\documentclass[11pt,a4paper]{article}

% ---------- Datos del documento ----------
\newcommand{\titulo}{Análisis PNI de las competencias del profesor y del tutor}
\newcommand{\tituloCorto}{Actividad 2}
\newcommand{\asignatura}{Organización y Gestión del Centro}
\newcommand{\autor}{Javier Gómez López}
\newcommand{\enunciado}{%
Una vez estudiadas las competencias del profesorado y de la figura del tutor en
el tema 5, realiza un análisis de las competencias del profesor y del tutor a
través de la estrategia \textbf{PNI: Positivo, Negativo, Interesante}.
\begin{itemize}[itemsep=0pt]
\item \textbf{P = positivo}: identifica los aspectos positivos de las dos figuras docentes y sus respectivas competencias.
\item \textbf{N = negativo}: identifica los aspectos negativos de las dos figuras docentes y sus respectivas competencias.
\item \textbf{I = interesante}: identifica los aspectos que encuentres interesantes respecto a las figuras docentes y sus competencias.
\end{itemize}
Redacta una conclusión que resulte de tu análisis (300 palabras máximo).
}
\newcommand{\master}{MÁSTER UNIVERSITARIO EN FORMACIÓN DEL PROFESORADO-ESPECIALIDAD EN MATEMÁTICAS -- Universidad Rey Juan Carlos}

% ---------- Paquetes ----------
\usepackage[utf8]{inputenc}
\usepackage[T1]{fontenc}
\usepackage[spanish,es-noshorthands]{babel}
\usepackage{lmodern}
\usepackage[margin=2.5cm,headheight=14pt]{geometry}
\usepackage[table]{xcolor}
\usepackage{titlesec}
\usepackage{fancyhdr}
\usepackage{enumitem}
\usepackage{parskip}
\usepackage{etoolbox}
\usepackage{booktabs}
\usepackage{longtable}
\usepackage{array}
\usepackage[most]{tcolorbox}
\usepackage[hidelinks]{hyperref}

% ---------- Estilo ----------
\definecolor{acento}{RGB}{176,0,32}   % rojo URJC
\definecolor{gris}{RGB}{90,90,90}
\definecolor{verde}{RGB}{46,125,50}
\definecolor{azul}{RGB}{21,101,192}
\definecolor{naranja}{RGB}{230,120,0}

\titleformat{\section}{\Large\bfseries\color{acento}}{\thesection.}{0.6em}{}
\titleformat{\subsection}{\large\bfseries\color{gris}}{\thesubsection.}{0.6em}{}

\pagestyle{fancy}
\fancyhf{}
\fancyhead[L]{\small\color{gris}\asignatura}
\fancyhead[R]{\small\color{gris}\tituloCorto}
\fancyfoot[C]{\small\thepage}
\renewcommand{\headrule}{{\color{acento}\hrule height 0.6pt}}

% Lista compacta para las celdas de la tabla PNI
\newlist{pni}{itemize}{1}
\setlist[pni]{label=--, leftmargin=1em, itemsep=2pt, topsep=2pt,
  parsep=0pt, before=\vspace*{-0.6\baselineskip}, after=\vspace*{0.5\baselineskip}}
\newcolumntype{C}{>{\raggedright\arraybackslash}p{0.305\linewidth}}
\newcommand{\figura}[1]{%
  \multicolumn{3}{l}{\cellcolor{acento!8}\textbf{\color{acento}#1}}\\*[2pt]}

\begin{document}

% ---------- Cabecera ----------
\thispagestyle{empty}
\begin{center}
  {\small\color{gris}\MakeUppercase{\master}}\\[1.2em]
  {\color{acento}\rule{\linewidth}{1.2pt}}\\[0.8em]
  {\Huge\bfseries \titulo}\\[0.6em]
  {\Large\color{gris} \asignatura}\\[0.6em]
  {\color{acento}\rule{\linewidth}{1.2pt}}\\[1em]
  {\large \autor}\\[0.3em]
  {\small\color{gris}\today}
\end{center}
\vspace{1.5em}

% ---------- Enunciado ----------
\ifdefempty{\enunciado}{}{%
  \begin{tcolorbox}[enhanced, breakable, colback=acento!4, colframe=acento,
      boxrule=0pt, leftrule=3pt, arc=0pt, title=Enunciado de la tarea,
      coltitle=acento, colbacktitle=acento!4, fonttitle=\bfseries]
    \enunciado
  \end{tcolorbox}
  \vspace{1em}
}

% ---------- Contenido ----------
\section{Las dos figuras en el tema 5}

El tema 5 no presenta al profesorado como un profesional aislado, sino como
miembro de una \textbf{estructura organizativa}. Los órganos de gobierno pueden
ser unipersonales (director, jefe de estudios, secretario) o colegiados
(Consejo Escolar y Claustro de profesores). Junto a ellos están los
\textbf{órganos de coordinación docente}: equipos docentes, áreas de
competencia, departamento de orientación, departamento de formación,
evaluación e innovación educativa, equipo técnico de coordinación pedagógica
(ETCP), tutor o tutora y departamentos de coordinación didáctica (T5,
diap.~2--3).

Las referencias del tipo \textbf{(T5, diap.~$n$)} remiten a la presentación
del tema 5. Los puntos sin referencia son \textbf{valoraciones propias}.

\begin{description}[leftmargin=0pt, style=nextline]
  \item[El profesor] Participa en varios órganos a la vez:
  \begin{itemize}[itemsep=0pt, topsep=2pt]
    \item En el \emph{departamento de coordinación didáctica} elabora la
    programación didáctica y hace su seguimiento, elabora, realiza y evalúa
    las pruebas específicas, evalúa la práctica docente y los procesos de
    enseñanza-aprendizaje y propone libros de texto y materiales (T5,
    diap.~12).
    \item En el \emph{equipo docente}, coordinado por el tutor, hace el
    seguimiento global del grupo, evalúa de forma colegiada, establece
    actuaciones para mejorar la convivencia y afronta los conflictos de manera
    coordinada (T5, diap.~4).
    \item A través de las \emph{áreas de competencia}, que agrupan los
    departamentos, trabaja en equipo las programaciones con una visión
    integrada y multidisciplinar (T5, diap.~5).
  \end{itemize}
  \item[El tutor o tutora] Cada grupo tiene uno, nombrado por el director a
  propuesta del jefe de estudios de entre los profesores que imparten clase al
  grupo. Desarrolla las actividades del Plan de Orientación y Acción Tutorial,
  mantiene una relación permanente con las familias, coordina las adaptaciones
  curriculares no significativas propuestas y elaboradas por el equipo docente
  y coordina el proceso de evaluación continua (T5, diap.~11). Para el plan de
  acción tutorial cuenta con el departamento de orientación, que colabora con
  el equipo directivo en su elaboración y asesora al profesorado en la
  atención a la diversidad y en la detección temprana de problemas de
  aprendizaje (T5, diap.~7).
\end{description}

\section{Análisis PNI}

La tabla separa las dos figuras para no mezclar lo que es propio de cada una,
aunque en la práctica casi todo tutor es, a la vez, profesor del grupo.

{\small
\renewcommand{\arraystretch}{1.15}
\setlength{\tabcolsep}{4pt}
\begin{longtable}{|C|C|C|}
\hline
\rowcolor{verde!12}
\multicolumn{1}{|c|}{\textbf{\color{verde}Positivo}} &
\multicolumn{1}{c|}{\textbf{\color{acento}Negativo}} &
\multicolumn{1}{c|}{\textbf{\color{azul}Interesante}} \\
\hline
\endhead
\hline
\endfoot
\figura{Profesor}
\begin{pni}
  \item \textbf{Autonomía profesional}: en el departamento elabora la
  programación, las pruebas y propone los libros y materiales (T5, diap.~12).
  \item \textbf{Trabajo colegiado}: la evaluación se hace de manera colegiada
  en el equipo docente (T5, diap.~4), lo que da una visión del alumno más
  completa que la de una sola materia.
  \item \textbf{Evaluación de la propia práctica}: el departamento de
  coordinación didáctica evalúa la práctica docente y los procesos de
  enseñanza-aprendizaje (T5, diap.~12).
\end{pni}
&
\begin{pni}
  \item \textbf{Acumulación de funciones}: enseñar, evaluar, atender a la
  diversidad, gestionar conflictos y coordinarse en varios órganos a la vez,
  con el mismo tiempo.
  \item \textbf{Riesgo de compartimentos estancos}: el departamento reúne al
  profesorado que imparte las mismas enseñanzas (T5, diap.~12), y la visión
  multidisciplinar depende de que funcionen las áreas de competencia.
\end{pni}
&
\begin{pni}
  \item Las <<competencias>> del profesor aparecen en el tema como
  \textbf{funciones dentro de órganos}: ser docente es una tarea colectiva.
  Un mismo profesor pertenece a la vez a un departamento, a un área de
  competencia, a varios equipos docentes y al Claustro.
  \item El departamento de formación, evaluación e innovación evalúa al
  claustro y el proceso de enseñanza, marca indicadores de calidad y propone
  planes de mejora (T5, diap.~8): una \textbf{evaluación entre iguales}
  dentro del propio centro.
\end{pni}
\\
\begin{pni}
  \item \textbf{Formación ajustada a necesidades}: el departamento de
  formación, evaluación e innovación diagnostica las necesidades formativas
  del profesorado y elabora los proyectos de formación (T5, diap.~8).
  \item \textbf{Papel en la convivencia}: el equipo docente establece
  actuaciones para mejorar la convivencia del grupo (T5, diap.~4).
  \item \textbf{Voz en la línea pedagógica}: los jefes de departamento forman
  parte del ETCP, que marca las pautas de las programaciones (T5, diap.~9).
\end{pni}
&
\begin{pni}
  \item La formación inicial suele ser muy disciplinar; competencias como la
  gestión de conflictos o la detección de problemas se aprenden
  <<sobre la marcha>>.
  \item La evaluación de la práctica docente puede quedarse en un trámite
  (memorias de fin de curso) si no se traduce en planes de mejora.
  \item Las programaciones del departamento pueden heredarse año tras año y
  frenar la innovación individual.
\end{pni}
&
\begin{pni}
  \item Ese mismo departamento investiga buenas prácticas y fomenta el
  \textbf{trabajo cooperativo entre departamentos} (T5, diap.~8): la
  innovación se plantea como tarea del centro, no solo individual.
  \item Las áreas de competencia buscan potenciar métodos que ayuden a
  alcanzar las \textbf{competencias} asignadas a cada área (T5, diap.~5): el trabajo
  del profesor se organiza en torno a ellas.
\end{pni}
\\ \hline
\noalign{\vspace{12pt}}
\hline
\figura{Tutor o tutora}
\begin{pni}
  \item \textbf{Adulto de referencia} para el grupo, elegido entre quienes le
  dan clase (T5, diap.~11): conoce al grupo de primera mano.
  \item \textbf{Visión global}: coordina la evaluación continua (T5,
  diap.~11) y al equipo docente, que hace el seguimiento global del grupo
  (T5, diap.~4). Hace de puente entre las materias.
  \item \textbf{Puente con las familias}: relación permanente con padres,
  madres o tutores legales (T5, diap.~11).
  \item \textbf{Orientación}: desarrolla las actividades del Plan de
  Orientación y Acción Tutorial (T5, diap.~11).
  \item \textbf{No trabaja solo}: el departamento de orientación colabora en
  el plan de acción tutorial y asesora al profesorado en la atención a la
  diversidad (T5, diap.~7).
  \item Coordina las adaptaciones curriculares no significativas (T5,
  diap.~11): es una pieza clave de la \textbf{inclusión}.
\end{pni}
&
\begin{pni}
  \item \textbf{Tiempo escaso} para todo lo que se le pide: la hora de
  tutoría, la atención a familias, las faltas, las actas y las
  comunicaciones.
  \item \textbf{Responsabilidad sin autoridad}: coordina al equipo docente
  (T5, diap.~4), pero sus compañeros no dependen de él, por lo que necesita
  su colaboración.
  \item Falta formación para \textbf{casos complejos} (ansiedad, absentismo,
  violencia, como en los casos del tema): debe saber detectar y derivar, pero
  no puede resolverlo todo.
  \item El nombramiento es por grupo y curso: si cambia cada año, el alumno
  puede perder su \textbf{referencia}.
\end{pni}
&
\begin{pni}
  \item La tutoría no es tarea de un especialista: el tutor sale del propio
  profesorado del grupo (T5, diap.~11). \textbf{Cualquier profesor puede ser
  tutor.}
  \item La acción tutorial es un \textbf{sistema}, no una persona: el tutor
  desarrolla un plan que elaboran el equipo directivo y el departamento de
  orientación (T5, diap.~7 y 11).
  \item El tutor suele ser el \textbf{primero en detectar} los problemas que
  afectan a la escolarización, como muestran los casos reales del tema (T5,
  diap.~13--14).
  \item El tutor es a la vez \textbf{igual y coordinador} de sus compañeros:
  es un liderazgo pedagógico que no viene del cargo, sino de la relación.
\end{pni}
\\ \hline
\figura{Ambas figuras}
\begin{pni}
  \item Se complementan: el profesor aporta el conocimiento de la materia y el
  tutor la visión de la persona.
\end{pni}
&
\begin{pni}
  \item Si la coordinación falla, el alumno recibe mensajes contradictorios
  de cada profesor.
\end{pni}
&
\begin{pni}
  \item Todos los órganos de coordinación del tema se apoyan en el
  \textbf{trabajo en equipo}: ninguna de las dos figuras funciona aislada.
\end{pni}
\\
\end{longtable}
}

\section{Conclusión}

% Máximo 300 palabras
El análisis muestra que las competencias del profesor y del tutor no se
entienden por separado, sino dentro de la organización del centro. Lo más
positivo es que el profesorado tiene autonomía para programar y evaluar su
materia y, a la vez, participa en las decisiones pedagógicas a través del
departamento, el área de competencia, el equipo docente y el Claustro. El
tutor añade la mirada sobre la persona: coordina la evaluación continua y al
equipo docente, desarrolla la acción tutorial y es el puente con las familias.

El principal punto débil es el mismo en las dos figuras: se les pide mucho
(atención a la diversidad, convivencia, coordinación, detección de
situaciones de riesgo) sin que crezcan en la misma medida el tiempo y la
formación específica. En el caso del tutor se suma una responsabilidad sin
autoridad formal sobre sus compañeros y un nombramiento anual que no siempre
garantiza la continuidad.

Lo más interesante es que la tutoría no es una tarea de especialistas:
cualquier profesor del grupo puede ser su tutor. Por eso, la acción tutorial
funciona como un sistema en el que el tutor, el departamento de orientación,
el equipo directivo y el equipo docente tienen que coordinarse. Si alguno
falla, el alumno lo nota.

Como futuro profesor de Matemáticas, saco dos ideas. La primera, que mi
formación disciplinar es necesaria pero no suficiente: tendré que formarme de
manera permanente en tutoría, convivencia e inclusión. La segunda, que el
trabajo en equipo no es un añadido, sino la forma normal de ejercer la
profesión. Un buen profesor de una materia que no se coordina con el resto del
equipo docente difícilmente podrá ser un buen tutor.

\vspace{1em}
{\small\color{gris}\textbf{Fuente.} Tema 5 de la asignatura: \emph{El
profesorado, los departamentos y los equipos educativos} (presentación).}

\end{document}
